% Options for packages loaded elsewhere
\PassOptionsToPackage{unicode}{hyperref}
\PassOptionsToPackage{hyphens}{url}
\documentclass[
  a4paper,
]{ltjsarticle}
\usepackage{xcolor}
\usepackage[margin=25mm]{geometry}
\usepackage{amsmath,amssymb}
\setcounter{secnumdepth}{-\maxdimen} % remove section numbering
\usepackage{iftex}
\ifPDFTeX
  \usepackage[T1]{fontenc}
  \usepackage[utf8]{inputenc}
  \usepackage{textcomp} % provide euro and other symbols
\else % if luatex or xetex
  \usepackage{unicode-math} % this also loads fontspec
  \defaultfontfeatures{Scale=MatchLowercase}
  \defaultfontfeatures[\rmfamily]{Ligatures=TeX,Scale=1}
\fi
\usepackage{lmodern}
\ifPDFTeX\else
  % xetex/luatex font selection
\fi
% Use upquote if available, for straight quotes in verbatim environments
\IfFileExists{upquote.sty}{\usepackage{upquote}}{}
\IfFileExists{microtype.sty}{% use microtype if available
  \usepackage[]{microtype}
  \UseMicrotypeSet[protrusion]{basicmath} % disable protrusion for tt fonts
}{}
\makeatletter
\@ifundefined{KOMAClassName}{% if non-KOMA class
  \IfFileExists{parskip.sty}{%
    \usepackage{parskip}
  }{% else
    \setlength{\parindent}{0pt}
    \setlength{\parskip}{6pt plus 2pt minus 1pt}}
}{% if KOMA class
  \KOMAoptions{parskip=half}}
\makeatother
\usepackage{longtable,booktabs,array}
\newcounter{none} % for unnumbered tables
\usepackage{calc} % for calculating minipage widths
% Correct order of tables after \paragraph or \subparagraph
\usepackage{etoolbox}
\makeatletter
\patchcmd\longtable{\par}{\if@noskipsec\mbox{}\fi\par}{}{}
\makeatother
% Allow footnotes in longtable head/foot
\IfFileExists{footnotehyper.sty}{\usepackage{footnotehyper}}{\usepackage{footnote}}
\makesavenoteenv{longtable}
\usepackage{graphicx}
\makeatletter
\newsavebox\pandoc@box
\newcommand*\pandocbounded[1]{% scales image to fit in text height/width
  \sbox\pandoc@box{#1}%
  \Gscale@div\@tempa{\textheight}{\dimexpr\ht\pandoc@box+\dp\pandoc@box\relax}%
  \Gscale@div\@tempb{\linewidth}{\wd\pandoc@box}%
  \ifdim\@tempb\p@<\@tempa\p@\let\@tempa\@tempb\fi% select the smaller of both
  \ifdim\@tempa\p@<\p@\scalebox{\@tempa}{\usebox\pandoc@box}%
  \else\usebox{\pandoc@box}%
  \fi%
}
% Set default figure placement to htbp
\def\fps@figure{htbp}
\makeatother
\setlength{\emergencystretch}{3em} % prevent overfull lines
\providecommand{\tightlist}{%
  \setlength{\itemsep}{0pt}\setlength{\parskip}{0pt}}
% Glyph map for lualatex (newunicodechar). Maps symbols that may lack font glyphs.
\usepackage{newunicodechar}
\usepackage{amssymb}
\newunicodechar{∎}{\ensuremath{\blacksquare}}
\newunicodechar{✓}{\ensuremath{\checkmark}}
\newunicodechar{✗}{\ensuremath{\times}}
\newunicodechar{⟹}{\ensuremath{\Longrightarrow}}
\newunicodechar{⟺}{\ensuremath{\Longleftrightarrow}}
\newunicodechar{↪}{\ensuremath{\hookrightarrow}}
\newunicodechar{⌊}{\ensuremath{\lfloor}}
\newunicodechar{⌋}{\ensuremath{\rfloor}}
\newunicodechar{≃}{\ensuremath{\simeq}}
\newunicodechar{≈}{\ensuremath{\approx}}
\newunicodechar{≤}{\ensuremath{\leq}}
\newunicodechar{≥}{\ensuremath{\geq}}
\newunicodechar{≠}{\ensuremath{\neq}}
\newunicodechar{→}{\ensuremath{\rightarrow}}
\newunicodechar{←}{\ensuremath{\leftarrow}}
\newunicodechar{↔}{\ensuremath{\leftrightarrow}}
\newunicodechar{⊥}{\ensuremath{\perp}}
\newunicodechar{√}{\ensuremath{\surd}}
\newunicodechar{∑}{\ensuremath{\sum}}
\newunicodechar{∏}{\ensuremath{\prod}}
\newunicodechar{∞}{\ensuremath{\infty}}
\newunicodechar{∈}{\ensuremath{\in}}
\newunicodechar{∀}{\ensuremath{\forall}}
\newunicodechar{∣}{\ensuremath{\mid}}
\newunicodechar{γ}{\ensuremath{\gamma}}
\newunicodechar{λ}{\ensuremath{\lambda}}
\newunicodechar{μ}{\ensuremath{\mu}}
\newunicodechar{ρ}{\ensuremath{\rho}}
\newunicodechar{σ}{\ensuremath{\sigma}}
\newunicodechar{φ}{\ensuremath{\varphi}}
\newunicodechar{ψ}{\ensuremath{\psi}}
\newunicodechar{θ}{\ensuremath{\theta}}
\newunicodechar{Δ}{\ensuremath{\Delta}}
\newunicodechar{Π}{\ensuremath{\Pi}}
\newunicodechar{Λ}{\ensuremath{\Lambda}}
\newunicodechar{τ}{\ensuremath{\tau}}
\newunicodechar{ε}{\ensuremath{\varepsilon}}
\newunicodechar{ω}{\ensuremath{\omega}}
\newunicodechar{π}{\ensuremath{\pi}}
\newunicodechar{ν}{\ensuremath{\nu}}
\newunicodechar{±}{\ensuremath{\pm}}
\newunicodechar{×}{\ensuremath{\times}}
\newunicodechar{·}{\ensuremath{\cdot}}
\usepackage{bookmark}
\IfFileExists{xurl.sty}{\usepackage{xurl}}{} % add URL line breaks if available
\urlstyle{same}
\hypersetup{
  pdftitle={The inflation curve is exponential amplification of the floor's linear instability mode------the largest eigenvalue of the floor Jacobian and its non-normality / time-varying corrections},
  pdfauthor={Noriaki Kihara (WF System Co., Ltd.), ORCID 0009-0004-6753-4020},
  hidelinks,
  pdfcreator={LaTeX via pandoc}}

\title{The inflation curve is exponential amplification of the floor's
linear instability mode------the largest eigenvalue of the floor
Jacobian and its non-normality / time-varying corrections}
\author{Noriaki Kihara (WF System Co., Ltd.), ORCID 0009-0004-6753-4020}
\date{2026-09-12}

\begin{document}
\maketitle

\textbf{Series:} Foundations of Self-Consistent Relational-Wave Closed
Systems (Paper 6 of 10)\\
\textbf{Author:} Noriaki Kihara (WF System Co., Ltd.)　\textbf{Date:}
2026-09-12\\
\textbf{Version DOI:} 10.5281/zenodo.22729096\\
\textbf{Concept DOI:} 10.5281/zenodo.22729095\\
\textbf{ORCID:} 0009-0004-6753-4020\\
\textbf{Zenodo:} https://zenodo.org/records/22729096\\

\begin{center}\rule{0.5\linewidth}{0.5pt}\end{center}

\subsection{Abstract}\label{abstract}

We derive analytically, as the \textbf{linear instability} of a tiny
seed, the inflation-like rapid expansion (exponential growth of the
out-of-floor-plane component \(H_\perp/H\)) reported in prior work
{[}1{]}. When the largest eigenvalue \(|\mu_{\max}|\) of the linearized
map of the floor (the Jacobian \(J\): the \(2M\times2M\) real matrix
obtained by linearizing \(z\to\exp(\Delta\tau K(\varphi))z\) at the
floor) exceeds unity, it gives
\(H_\perp/H(t)\approx\text{seed}^2\,|\mu_{\max}|^{2t}\), ramp slope
\(2\log_{10}|\mu_{\max}|\), and onset
\(\approx\ln(1/H_{\perp0})/(2\ln|\mu_{\max}|)=\ln(1/\text{seed})/\ln|\mu_{\max}|\)
(with \(H_{\perp0}=\text{seed}^2\)). The \(H_\perp\) slope obtained by
evolving a tiny \(\delta\) under the frozen floor \(J\) matches
\(2\log_{10}|\mu_{\max}|\) exactly (\(N=6{:}0.2935\) vs analytic
\(0.2937\)). \(|\mu_{\max}|\) decreases as \(N\) increases
(\(1.402\to1.269\to1.185\) for \(N=6,8,10\)) = inflation becomes slower
and the onset lengthens (measured onset is \(64,87,133\) steps for
\(N=6,8,10\)). The measured ramp slope is consistently \(18\)--\(31\%\)
steeper than the floor value, and this stems from the fact that the
floor \(J\) is a strongly \textbf{non-normal matrix} (non-normality
\(0.53\)--\(0.76\)) and that the real ramp is the amplification of the
\textbf{time-varying, non-commuting time-ordered product}
\(\prod J(\varphi_t)\) accompanying the gradual change of \(K(\varphi)\)
(joint spectral radius \(\ge\) individual spectral radius)
(\textbf{directly confirmed in follow-up test A}: the slope of the
time-varying product matches the measured value and exceeds the
frozen-floor value by 18--31\%). The floor eigenvalue is its lower bound
(initial / frozen rate). Classification: derivable (form, mechanism,
\(N\) dependence, onset) + direct confirmation of the non-normality
correction (follow-up test A; only the closed form is open).

\begin{center}\rule{0.5\linewidth}{0.5pt}\end{center}

\subsection{1. Problem}\label{problem}

In Papers 1 and 2 the floor was established as an unstable relative
equilibrium. ``If the dynamics of the wave has been solved, then the
inflation curve too should be derivable analytically from a tiny
seed''------this paper shows this quantitatively, as the exponential
amplification of the linear instability mode of the floor Jacobian.

\subsection{2. The linear-instability prediction
formulas}\label{the-linear-instability-prediction-formulas}

We construct by numerical differentiation the real Jacobian
\(J\in\mathbb R^{2M\times2M}\) obtained by linearizing the map around
the floor \(z_0\) (embedding the state as \([\Re z;\Im z]\)), and
compute its eigenvalues \(\mu\). If the floor is unstable,
\(|\mu_{\max}|>1\). A tiny seed \(\delta\) is amplified along the
unstable eigendirection, and the fraction of the component orthogonal to
the floor plane is

\[\frac{H_\perp}{H}(t)\approx(\text{seed})^2\,|\mu_{\max}|^{2t},\qquad
\text{ramp slope}=\frac{d\log_{10}(H_\perp/H)}{dt}=2\log_{10}|\mu_{\max}|,\qquad
\text{onset}\approx\frac{\ln(1/H_{\perp0})}{2\ln|\mu_{\max}|}=\frac{\ln(1/\text{seed})}{\ln|\mu_{\max}|}.\]

(The onset is the number of steps to reach \(H_\perp/H\to1\).
\(H_{\perp0}=\text{seed}^2\). The old draft's
\(\ln(1/\text{seed})/(2\ln|\mu|)\) was a coefficient error that produced
a self-contradiction---about half of the measured onset---so it has been
corrected.)

\subsection{\texorpdfstring{3. Verification
(\(N=6,8,10\))}{3. Verification (N=6,8,10)}}\label{verification-n6810}

{\def\LTcaptype{none} % do not increment counter
\begin{longtable}[]{@{}
  >{\raggedright\arraybackslash}p{(\linewidth - 18\tabcolsep) * \real{0.1000}}
  >{\raggedright\arraybackslash}p{(\linewidth - 18\tabcolsep) * \real{0.1000}}
  >{\raggedright\arraybackslash}p{(\linewidth - 18\tabcolsep) * \real{0.1000}}
  >{\raggedright\arraybackslash}p{(\linewidth - 18\tabcolsep) * \real{0.1000}}
  >{\raggedright\arraybackslash}p{(\linewidth - 18\tabcolsep) * \real{0.1000}}
  >{\raggedright\arraybackslash}p{(\linewidth - 18\tabcolsep) * \real{0.1000}}
  >{\raggedright\arraybackslash}p{(\linewidth - 18\tabcolsep) * \real{0.1000}}
  >{\raggedright\arraybackslash}p{(\linewidth - 18\tabcolsep) * \real{0.1000}}
  >{\raggedright\arraybackslash}p{(\linewidth - 18\tabcolsep) * \real{0.1000}}
  >{\raggedright\arraybackslash}p{(\linewidth - 18\tabcolsep) * \real{0.1000}}@{}}
\toprule\noalign{}
\begin{minipage}[b]{\linewidth}\raggedright
\(N\)
\end{minipage} & \begin{minipage}[b]{\linewidth}\raggedright
\(|\mu_{\max}|\)
\end{minipage} & \begin{minipage}[b]{\linewidth}\raggedright
\(\lambda=\ln|\mu|\)/step
\end{minipage} & \begin{minipage}[b]{\linewidth}\raggedright
number of unstable modes
\end{minipage} & \begin{minipage}[b]{\linewidth}\raggedright
analytic slope \(2\log_{10}|\mu|\)
\end{minipage} & \begin{minipage}[b]{\linewidth}\raggedright
\(\delta\)-evolution slope under frozen floor \(J\)
\end{minipage} & \begin{minipage}[b]{\linewidth}\raggedright
measured ramp slope
\end{minipage} & \begin{minipage}[b]{\linewidth}\raggedright
measured/floor-analytic
\end{minipage} & \begin{minipage}[b]{\linewidth}\raggedright
floor \(J\) non-normality
\end{minipage} & \begin{minipage}[b]{\linewidth}\raggedright
onset analytic/measured
\end{minipage} \\
\midrule\noalign{}
\endhead
\bottomrule\noalign{}
\endlastfoot
6 & 1.40234 & 0.3381 & 5 & 0.2937 & 0.2935 & 0.3460 & 1.178 & 0.761 & 86
/ 64 \\
8 & 1.26908 & 0.2383 & 6 & 0.2070 & 0.1988 & 0.2702 & 1.306 & 0.641 &
125 / 87 \\
10 & 1.18513 & 0.1699 & 8 & 0.1475 & 0.1271 & 0.1814 & 1.230 & 0.530 &
180 / 133 \\
\end{longtable}
}

\begin{itemize}
\tightlist
\item
  \textbf{The linear-instability formula is exactly correct}: the
  \(H_\perp\) slope from evolving a tiny \(\delta\) under the frozen
  floor \(J\) matches the analytic \(2\log_{10}|\mu_{\max}|\)
  (\(N=6{:}0.2935\) vs \(0.2937\)).
\item
  \textbf{The \(N\) dependence is reproduced}: \(|\mu_{\max}|\)
  decreases as \(N\) increases (\(1.402\to1.269\to1.185\)) = inflation
  becomes slower and the onset lengthens (measured onset is
  \(64,87,133\) steps for \(N=6,8,10\), analytic \(86,125,180\) steps).
  The closed-form \(N\)-dependence law of the onset is left as future
  work.
\item
  \textbf{onset} \(\approx\ln(1/\text{seed})/\ln|\mu|\) matches in order
  of magnitude and scale (analytic \(86/125/180\) vs measured
  \(64/87/133\)).
\end{itemize}

\begin{figure}
\centering
\pandocbounded{\includegraphics[keepaspectratio,alt={Figure 1: Inflation = linear instability. Left = N dependence of the floor Jacobian \textbar\textbackslash mu\_\{\textbackslash max\}\textbar{} and the onset; center = analytic ramp slope 2\textbackslash log\_\{10\}\textbar\textbackslash mu\textbar{} (floor) vs measured (non-normality gap); right = the measured H\_\textbackslash perp/H ramp for N=6 with the analytic slope overlaid.}]{figs/p6_inflation_linear_instability_en_fig1.png}}
\caption{Figure 1: Inflation = linear instability. Left = \(N\)
dependence of the floor Jacobian \(|\mu_{\max}|\) and the onset; center
= analytic ramp slope \(2\log_{10}|\mu|\) (floor) vs measured
(non-normality gap); right = the measured \(H_\perp/H\) ramp for \(N=6\)
with the analytic slope overlaid.}
\end{figure}

\textbf{Figure 1} Inflation = the floor's linear instability (\(N\)
dependence, slope, overlay).

\subsection{4. Non-normality correction (the amount by which the
measurement exceeds the floor
value)}\label{non-normality-correction-the-amount-by-which-the-measurement-exceeds-the-floor-value}

The measured ramp slope is consistently \(18\)--\(31\%\) steeper than
the floor value (measured/floor-analytic \(=1.18\)--\(1.31\)).

\begin{itemize}
\tightlist
\item
  The floor \(J\) is a strongly \textbf{non-normal matrix}
  (non-normality
  \(\|JJ^{\mathsf T}-J^{\mathsf T}J\|/\|J^{\mathsf T}J\|=0.53\)--\(0.76\)).
\item
  The spectral radius when \(J\) is frozen at each point is nearly
  constant along the ramp (\(1.402\) for \(N=6\)), but the real ramp
  \textbf{persistently} exceeds it. Frozen non-normality only provides
  transient amplification; the persistent part is carried by the
  amplification of the \textbf{time-varying, non-commuting time-ordered
  product of Jacobians} \(\prod_t J(\varphi_t)\) (joint spectral radius
  \(\ge\) individual spectral radius) as \(K(\varphi)\) changes
  gradually. \textbf{This mechanism has been established by direct
  verification (follow-up test A)}: the \(H_\perp\) slope from evolving
  \(\delta\) under the time-varying Jacobian product \(\prod_t J(z_t)\)
  on the real trajectory matches the measured ramp
  (time-varying/measured = 0.999 / 1.004 / 1.019 for \(N=6,8,10\)), and
  exceeds the frozen floor \(J_0\) by 18--31\% (time-varying/floor =
  1.177 / 1.311 / 1.253 = quantitatively matching the measured excess).
  The floor eigenvalue is its lower bound (initial / frozen rate).
  Therefore the fact that the analytic onset (based on the floor
  \(|\mu_{\max}|\)) in §3 consistently exceeds the measured value
  (\(86/125/180\) vs measured \(64/87/133\)) is likewise a consequence
  of the same non-normality mechanism------because the floor eigenvalue
  is a lower bound on the growth rate, the floor-based onset gives an
  upper bound, and the real ramp is steeper (slope ratio
  \(1.18\)--\(1.31\)) and reaches the threshold faster. The onset
  overestimation and the non-normality excess are two sides of the same
  mechanism.
\end{itemize}

\textbf{Deeper reading (finite-time Lyapunov rate of the tangent
cocycle).} Near the floor, \(\delta_{t+1}=J_0\delta_t\) and the rate is
\(\ln\rho(J_0)=\ln|\mu_{\max}|\) (the initial growth rate near the
floor). Once away from the floor, on the real trajectory
\(\delta_{t+1}=J(z_t)\delta_t\), i.e.~\(\delta_t=\Phi(t,0)\delta_0\),
\(\Phi(t,0)=J(z_{t-1})\cdots J(z_0)\) (the time-ordered product of
tangent maps = the \textbf{tangent cocycle}). Hence \textbf{the real
inflation rate is the finite-time Lyapunov rate of the tangent cocycle
\(\Phi(t,0)=\prod_t J(z_t)\) on the real trajectory}

\[\gamma_t=\frac1t\ln\frac{\|P_\perp\Phi(t,0)\delta_0\|}{\|P_\perp\delta_0\|},\qquad \text{slope of }\log_{10}(H_\perp/H)=\frac{2\gamma_t}{\ln10}\]

(\(P_\perp\) = orthogonal projection onto the floor plane). Since the
\(J_t\) are non-normal (\(0.53\)--\(0.76\)) and mutually non-commuting,
\(\Phi(t,0)\neq J_0^t\), and the direction of maximal amplification is
further amplified by the next \(J_t\) while rotating under time
evolution (unstable eigenmode → rotating amplified direction →
time-ordered exponential growth). Follow-up test A confirmed this
\(\gamma_t\) as matching the measured ramp (time-varying/measured 0.999
/ 1.004 / 1.019). \(\mu_{\max}\) (the local instability of the floor)
and \(\Phi(t,0)\) (the tangent dynamics of the real inflation) are a
\textbf{two-stage dynamics}, and \(\mu_{\max}\) is its lower bound
(\(t\to0^+\) / frozen limit).

\subsection{5. Conclusion of the two-stage
structure}\label{conclusion-of-the-two-stage-structure}

\begin{itemize}
\tightlist
\item
  \textbf{Derivable (form, mechanism, \(N\) dependence, onset)}: the
  inflation curve is the exponential amplification of the floor's linear
  instability mode. The main predictions (slope \(2\lambda/\ln10\),
  onset, \(N\) dependence) are given by the largest eigenvalue of the
  floor Jacobian \(\lambda=\ln|\mu_{\max}|\).
\item
  \textbf{Unification with Paper 5 (two faces of the same
  \(\mu_{\max}\))}: this \(\lambda=\ln|\mu_{\max}|\) is precisely the
  summit curvature of the Mexican-hat effective potential of Paper 5
  (\(V''(0)=-\lambda\), \(\lambda=\ln|\mu_{\max}|\)). That is,
  \textbf{the instability of SSB (the negative curvature of the summit
  that drives the fall into the valley) and the exponential growth rate
  of inflation are tied together by the same floor Jacobian eigenvalue
  \(\mu_{\max}\)}. Paper 5's static SSB picture (spontaneous selection
  from an unstable symmetric maximum) and this paper's dynamic
  amplification (exponential rapid expansion) are the static and dynamic
  aspects of a single \(\mu_{\max}\). As \(N\) increases and
  \(\mu_{\max}\to1\), the summit flattens (curvature decreases), and the
  fall = inflation also becomes slower (onset lengthens). That the
  growth rate does not depend on precision or seed depth (= it is
  intrinsic to the dynamics) has already been established in an
  experiment that verified the floor of the seed at 100-digit
  precision------even at IC100 (closure \(\sim10^{-102}\), floor
  \(\sim10^{-200}\)), \(\gamma_\tau\) agrees with the 64-bit run to 8
  digits (\(N=8\)), and it fires across about 179 digits from the floor
  under a single exponential law
  (\texttt{precision\_seed\_dynamics\_isolation\_100digit\_20260902},
  \(N=7,8\)).
\item
  \textbf{Refinement (non-normality correction)}: the exact slope is the
  amplification of the time-ordered product of the time-varying
  non-normal \(J\), adding \(\sim18\)--\(31\%\) on top of the floor
  eigenvalue (\textbf{directly confirmed in follow-up test A}: the slope
  of the time-varying product matches the measured ramp and exceeds the
  frozen floor by 18--31\%). Only the closed form remains underived.
\end{itemize}

\textbf{Unified chain} (a single line from ignition to landing):

\[\rho(J_0)>1\ \Rightarrow\ \text{ignition qualification (Paper 2)}\ \to\ \ln\rho(J_0)\ \Rightarrow\ \text{initial growth rate near the floor}\ \to\ \gamma_{\rm FT}\big[\textstyle\prod_t J(z_t)\big]\ \Rightarrow\ \text{real ramp growth rate (follow-up test A)}\ \to\ \frac{\ln(1/\text{seed})}{\gamma}\ \Rightarrow\ \text{onset}\ \to\ \text{nonlinear saturation}\ \Rightarrow\ \text{degenerate vacuum (Paper 5)}.\]

Viewing the same unstable floor by ``how it leaves the floor'' gives
inflation, and by ``where it lands'' gives SSB------the two are not
distinct mechanisms but two views of a single unstable floor.

This is a quantitative description of rolling down from the Mexican-hat
summit (the unstable relative equilibrium of the 90° skeleton), and the
landing (the degenerate vacuum, \(\Delta=-2\pi/N\)) is treated by Paper
5, while the geometry of tilting away from the floor plane is treated by
Paper 7.

\subsection{6. Claim classification / open
items}\label{claim-classification-open-items}

\begin{itemize}
\tightlist
\item
  Classification: derived consequence (the linear-instability formula) +
  measurement of the non-normality correction (the refinement is a
  hypothesis). Judgment: retained.
\item
  Open: the \textbf{closed form} of the amplification rate of the
  non-normal time-varying product is underived (the mechanism is
  established by follow-up test A; only the analytic expression of the
  amplification rate remains for future work). The floor Jacobian for
  \(N=40\) (\(2M=1560\) dimensions) is heavy and uncomputed------for
  small \(N\) (6, 8, 10) the two-stage structure is established.
\end{itemize}

\subsection{Related work (structural correspondence, not the source of
derivation)}\label{related-work-structural-correspondence-not-the-source-of-derivation}

\begin{itemize}
\tightlist
\item
  Transient amplification / pseudospectra of non-normal matrices
  (Trefethen \& Embree 2005): the amount by which the ramp exceeds the
  floor eigenvalue is the non-normal transient growth / amplification of
  the time-ordered product.
\item
  Self-consistent inflation v2 {[}1{]}: this paper is the analytic
  derivation (refinement) of its rapid-expansion mechanism.
\end{itemize}

\subsection{References}\label{references}

\begin{enumerate}
\def\labelenumi{\arabic{enumi}.}
\tightlist
\item
  Noriaki Kihara, ``The mechanism of inflation-like rapid expansion in
  self-consistent relational-wave closed systems,'' Version DOI
  10.5281/zenodo.22176949 / Concept DOI 10.5281/zenodo.22112008.
\item
  L. N. Trefethen and M. Embree, \emph{Spectra and Pseudospectra},
  Princeton (2005).
\end{enumerate}

\subsection{Reproduction}\label{reproduction}

The verification scripts, figures, and reports are in
\texttt{位相ロック全N確認\_相対平衡\_20260912/インフレ曲線\_線形不安定解析\_20260912/}
(\texttt{analyze\_inflation\_linear\_instability\_20260912.py} =
verifications A--E {[}floor \(J\) eigenvalues / non-normality /
frozen-floor \(\delta\) evolution / measured ramp / onset{]},
\texttt{make\_figure\_inflation.py},
\texttt{fig\_inflation\_linear\_instability.png}, \texttt{REPORT.md},
\texttt{解析ログ\_20260912.log}, \texttt{SHA256SUMS.txt},
\texttt{run\_all.sh}). The \textbf{direct confirmation of the mechanism
of the non-normality correction (follow-up test A)} is in
\texttt{位相ロック全N確認\_相対平衡\_20260912/ランプ直接Lyapunov\_非正規時変積\_20260912/}
(\texttt{analyze\_ramp\_lyapunov\_20260912.py} = comparison of the
\(H_\perp\) slopes of frozen floor / time-varying product / measurement,
\texttt{results/ramp\_lyapunov\_summary.csv}, \texttt{README.md} /
\texttt{REPORT.md} / \texttt{SHA256SUMS.txt} / \texttt{run\_all.sh}).
The input is the existing den=N npz (500 steps, make\_parent floor). No
runs or physics were changed. The dynamical map is identical to the
canonical \texttt{run\_N3\_N40\_stage123\_v1.py} (SHA256
\texttt{1abf2353fee2e4f56f05e7a6f149fd086885136beb61ab571b48a56b09691567}).

\end{document}
